% =============================================================================
\section{Results}
\label{sec:results}
% =============================================================================

We evaluate on 31 datasets across five domains, each attacked with three MIA
methods (LiRA, shadow, loss-threshold) under three model families (MLP,
XGBoost, Random Forest), for 279 (dataset, attack, model) configurations.

\subsection{Experimental Ground Truth}

We define the ground-truth risk of a dataset as the mean MIA AUC over the
nine attack-model configurations, $\mathrm{Risk}(D)$.
We use the mean rather than the per-dataset maximum because the maximum
collapses onto the Random Forest for every dataset. Random Forest with no
depth limit memorizes its training set, so its AUC saturates near $1.0$ and
erases all contrast between datasets. The mean preserves that contrast.

Aggregated over the grid, LiRA is the strongest attack (mean AUC $0.71$)
ahead of loss-threshold and shadow ($0.61$ each), and Random Forest is the
most vulnerable model (mean AUC $0.75$) ahead of XGBoost ($0.62$) and MLP
($0.56$). $\mathrm{Risk}(D)$ ranges almost two-fold across datasets, from
$0.92$ (Texas100) down to $0.51$ (Gowalla, Mushroom), confirming that the
choice of dataset alone moves privacy risk substantially.

\subsection{DPRI Ranks Privacy Risk}

We predict $\mathrm{Risk}(D)$ from the DPRI features using ridge regression
under leave-one-dataset-out cross-validation. Because the geometric features
are heavily right-skewed (local density spans $0.03$ to $35$), we
rank-transform features within each fold and report Spearman rank
correlation, which is the relevant metric for a pre-training risk index that
only needs to order datasets.

DPRI ranks datasets by risk with a Spearman correlation of $0.61$ under an
unbiased nested cross-validation that selects the feature subset on training
folds only (Table~\ref{tab:regression}). The fixed-feature LOO estimate is
$0.52$ with all features and $0.61$ with the geometric-core-plus-dimension
subset. Random Forest regression does worse ($\rho=0.36$), overfitting the
small sample. The prediction is moderate, not precise. Coefficients of
determination are not an informative lens here, since the target is a
ranking and the features are rank-transformed.

\begin{table}[t]
\centering
\caption{DPRI ranks $\mathrm{Risk}(D)$ across 31 datasets (Spearman $\rho$,
         leave-one-dataset-out). The nested-CV number is the unbiased
         headline.}
\label{tab:regression}
\small
\begin{tabular}{lc}
\toprule
\textbf{Estimator} & \textbf{Spearman} $\rho$ \\
\midrule
Nested-CV (unbiased, headline)        & $0.61$ \\
Ridge, all six features (fixed)       & $0.52$ \\
Ridge, geometric core + dimension     & $0.61$ \\
Random Forest, all features           & $0.36$ \\
\bottomrule
\end{tabular}
\end{table}

Table~\ref{tab:dpri_features} illustrates the relationship on representative
datasets. High-risk datasets tend to have high sample uniqueness and low
density (Texas100, Purchase100), and low-risk datasets the reverse (Gowalla,
Mushroom). The relationship is a tendency rather than a law: Creditg has
high uniqueness yet only moderate risk, and Mushroom has moderate uniqueness
yet the lowest risk. These exceptions are why the correlation is $0.61$
rather than $1.0$.

\begin{table}[t]
\centering
\caption{DPRI features and $\mathrm{Risk}(D)$ for representative datasets,
         ordered by risk. $u$: sample uniqueness, $\rho$: local density,
         $S$: class separability, $d$: dimensionality.}
\label{tab:dpri_features}
\small
\setlength{\tabcolsep}{4pt}
\begin{tabular}{lrrrrr}
\toprule
\textbf{Dataset} & $\bar{u}$ & $\bar{\rho}$ & $S$ & $d$ & \textbf{Risk} \\
\midrule
Texas100     & $55.7$ & $0.03$ & $-0.25$ & $6169$ & $0.92$ \\
Purchase100  & $28.0$ & $0.04$ & $-0.02$ &  $600$ & $0.85$ \\
Creditg      &  $6.9$ & $0.15$ & $ 0.02$ &   $74$ & $0.77$ \\
Adult        &  $1.0$ & $1.55$ & $-0.01$ &   $14$ & $0.70$ \\
MNIST        & $16.2$ & $0.08$ & $-0.05$ &  $784$ & $0.66$ \\
Magic        &  $0.8$ & $1.49$ & $ 0.12$ &   $10$ & $0.63$ \\
Movielens    &  $0.2$ & $6.10$ & $ 0.12$ &    $5$ & $0.53$ \\
Gowalla      &  $0.1$ & $34.9$ & $ 0.09$ &    $6$ & $0.51$ \\
Mushroom     &  $3.8$ & $0.28$ & $ 0.09$ &  $138$ & $0.51$ \\
\bottomrule
\end{tabular}
\end{table}

\begin{figure}[t]
  \centering
  \dprifig{figures/regression_plot.pdf}{\linewidth}
  \caption{Predicted vs.\ observed $\mathrm{Risk}(D)$ under
           leave-one-dataset-out cross-validation. Each point is one of the
           31 datasets.}
  \label{fig:regression}
\end{figure}

\subsection{Which Factors Carry the Signal}
\label{sec:single_factor}

We assess each feature by its standalone Spearman correlation with
$\mathrm{Risk}(D)$ (Table~\ref{tab:single_factor}). A drop-one ablation is
not meaningful here, because the features are collinear and the linear
baseline is weak, so removing one feature lets a correlated feature absorb
its signal. Standalone rank correlation is robust to both.

Three geometric features are individually significant: local density,
class separability, and sample uniqueness (all $p < 0.005$). Outlier score,
feature entropy, and dimensionality are not individually significant at
$n=31$, though Random Forest assigns them nonzero importance, suggesting
they contribute weak or nonlinear signal. The headline result is that
no single statistic determines risk, and that the significant factors are
precisely the geometric ones our theory predicts.

\begin{table}[t]
\centering
\caption{Standalone ranking power of each DPRI feature (Spearman vs.\
         $\mathrm{Risk}(D)$, $n=31$).}
\label{tab:single_factor}
\small
\begin{tabular}{lcc}
\toprule
\textbf{Feature} & \textbf{Spearman} $\rho$ & \textbf{$p$-value} \\
\midrule
Local density $\bar{\rho}$      & $-0.565$ & $0.001$ \\
Class separability $S$          & $-0.543$ & $0.002$ \\
Sample uniqueness $\bar{u}$     & $+0.506$ & $0.004$ \\
Dimensionality $\log d$         & $+0.294$ & $0.108$ \\
Outlier score $\bar{o}$         & $+0.292$ & $0.111$ \\
Feature entropy $H$             & $-0.084$ & $0.654$ \\
\bottomrule
\end{tabular}
\end{table}

\subsection{Finding 1: Dataset Structure Governs Risk Among Realistic Models}
\label{sec:finding1}

We decompose the variation in MIA AUC by the spread of group-mean AUC across
each factor. Across all nine configurations, dataset choice accounts for
$36\%$ of the variation, model architecture for $39\%$, and the attack
method for $24\%$. Taken at face value, the model appears to rival the
dataset.

This parity is almost entirely an artifact of a single model. The Random
Forest, trained without a depth limit, memorizes its training set and
saturates near AUC $1.0$ on nearly every dataset. Its mean AUC is $0.75$,
against $0.56$ for the MLP and $0.62$ for XGBoost. This is an adversarial
worst case, not a configuration that practitioners deploy. Restricting the
decomposition to the regularized models that are actually used, MLP and
XGBoost, dataset structure rises to $49.5\%$ of the variation, twice the
model share ($24.4\%$) and well ahead of attack ($26.1\%$).

We therefore read Finding~1 as follows. Across realistic models, dataset
structure governs membership inference risk more than the choice of model,
and far more than the choice of attack. A deliberately overfit model can
erase that margin, which is itself a caution. Reporting a single attack
number on a single dataset conflates three distinct sources of variation,
and the dataset is not the incidental one.

\subsection{Finding 2: Benchmarks Are Geometric Outliers}
\label{sec:finding2}

The two most widely used MIA benchmarks, Purchase100 and Texas100, are
geometric outliers among the 31 datasets. Their sample uniqueness ($28.0$
and $55.7$) is the highest of all datasets we study, far above the median
of roughly $1.5$, and their risk ($0.85$ and $0.92$) is likewise the top
two. A model evaluated only on these benchmarks is being attacked in an
unrepresentatively high-risk regime.

We temper this in two ways. First, the effect is not exclusive to
benchmarks: high-dimensional image datasets such as MNIST and Fashion-MNIST
also have elevated uniqueness, because high dimensionality sparsifies the
data. Second, with only two benchmark datasets a formal distributional test
has little power. The claim we make is therefore concrete and limited:
on the single most important geometric axis, sample uniqueness, the two
canonical MIA benchmarks sit at the extreme of the entire 31-dataset range.
Corollary~\ref{cor:benchmark} explains why subsampled, class-balanced
benchmarks are expected to have inflated uniqueness from first principles.
